\section{Merge K Sorted Lists}
\label{sec:heap-merge-k-lists}

\problemstatement{We are given $k$ linked lists, each individually sorted
in increasing order. Merge all of them into a single sorted linked list
and return it.}

\begin{patternbox}
\emph{``Merge $k$ sorted ...''} is one of the cleanest heap triggers in
this chapter: merging \emph{two} sorted lists is a familiar $O(n)$
two-pointer operation, but merging $k$ of them at once naturally wants a
data structure that can always hand you the smallest \emph{current head}
among all $k$ lists in $O(\log k)$ time -- exactly a min-heap holding one
candidate node from each list.
\end{patternbox}

\subsection*{Understanding the problem carefully}

At any point during the merge, the next node to append to the output is
the smallest value among the \emph{current heads} of the $k$ lists that
still have remaining nodes (every list is individually sorted, so only its
current head -- not any later node -- can possibly be the next-smallest
value globally). A naive approach would scan all $k$ current heads every
time to find the minimum, costing $O(k)$ per output node and $O(nk)$
overall (where $n$ is the total number of nodes). A min-heap reduces the
``find the minimum among $k$ candidates'' step to $O(\log k)$.

\begin{ideabox}[Greedy Choice, Powered by a Min-Heap]
Initialize a min-heap with the head node of each of the $k$ lists (if
non-null). Repeatedly: pop the smallest node from the heap, append it to
the output list, and if that node has a successor within its own original
list, push that successor onto the heap. Continue until the heap is
empty.
\end{ideabox}

\subsection*{Why popping the heap's minimum is always the correct next output node}

The heap, at every step, holds exactly one ``current head'' candidate per
list that still has unconsumed nodes -- and because each list is
individually sorted, the true global minimum of all \emph{remaining}
nodes across all lists must be one of these $k$ current heads (a node
deeper inside a list cannot be smaller than that same list's own current
head). So the heap's minimum at any moment is guaranteed to be the correct
next node for the merged output, and pushing that node's successor
immediately after popping it maintains the invariant that the heap always
holds the current head of every still-nonempty list.

\subsection*{Algorithm}

\begin{enumerate}
  \item Initialize a min-heap (ordered by node value) with the head of
        each of the $k$ lists, skipping any empty lists.
  \item Initialize an empty output list (with a dummy head for easy
        appending).
  \item While the heap is not empty:
  \begin{enumerate}
    \item Pop the node with the smallest value; append it to the output
          list.
    \item If this node has a \texttt{next} pointer (within its original
          list): push that next node onto the heap.
  \end{enumerate}
  \item Return the output list (skipping the dummy head).
\end{enumerate}

\begin{algorithm}[H]
\caption{Merge K Sorted Lists}
\begin{algorithmic}[1]
\Procedure{MergeKLists}{\textit{lists}}
  \State $\textit{heap} \gets$ empty min-heap (by node value)
  \For{each list $L$ in \textit{lists}}
    \If{$L$ is not empty} push $L$'s head onto \textit{heap} \EndIf
  \EndFor
  \State $\textit{dummy} \gets$ new empty node; $\textit{tail} \gets \textit{dummy}$
  \While{\textit{heap} is not empty}
    \State $\textit{node} \gets$ pop minimum from \textit{heap}
    \State $\textit{tail}.\textit{next} \gets \textit{node}$; $\textit{tail} \gets \textit{node}$
    \If{$\textit{node}.\textit{next} \ne \text{null}$}
      \State push $\textit{node}.\textit{next}$ onto \textit{heap}
    \EndIf
  \EndWhile
  \State \Return $\textit{dummy}.\textit{next}$
\EndProcedure
\end{algorithmic}
\end{algorithm}

\subsection*{Dry run}

\dryrun{Take three lists: $L_1 = 1\to4\to5$, $L_2=1\to3\to4$,
$L_3=2\to6$.}

\begin{itemize}
  \item Initialize heap with the three heads: $\{1_{(L_1)}, 1_{(L_2)},
        2_{(L_3)}\}$.
  \item Pop a $1$ (say from $L_1$): output $=1$. Push $L_1$'s next, $4$.
        Heap $=\{1_{(L_2)}, 2_{(L_3)}, 4_{(L_1)}\}$.
  \item Pop $1$ (from $L_2$): output $=1,1$. Push $L_2$'s next, $3$.
        Heap $=\{2_{(L_3)}, 3_{(L_2)}, 4_{(L_1)}\}$.
  \item Pop $2$ (from $L_3$): output $=1,1,2$. Push $L_3$'s next, $6$.
        Heap $=\{3_{(L_2)}, 4_{(L_1)}, 6_{(L_3)}\}$.
  \item Pop $3$ (from $L_2$): output $=1,1,2,3$. Push $L_2$'s next, $4$.
        Heap $=\{4_{(L_1)}, 4_{(L_2)}, 6_{(L_3)}\}$.
  \item Pop a $4$ (say from $L_1$): output $=1,1,2,3,4$. Push $L_1$'s
        next, $5$. Heap $=\{4_{(L_2)}, 5_{(L_1)}, 6_{(L_3)}\}$.
  \item Pop $4$ (from $L_2$): output $=1,1,2,3,4,4$. $L_2$ exhausted, push
        nothing. Heap $=\{5_{(L_1)}, 6_{(L_3)}\}$.
  \item Pop $5$: output $=\dots,5$. $L_1$ exhausted. Heap $=\{6_{(L_3)}\}$.
  \item Pop $6$: output $=\dots,6$. Heap empty.
\end{itemize}

Final merged list: $1,1,2,3,4,4,5,6$.

\subsection*{C++ implementation}

\begin{lstlisting}
#include <bits/stdc++.h>
using namespace std;

struct ListNode {
    int val;
    ListNode* next;
    ListNode(int v) : val(v), next(nullptr) {}
};

struct Compare {
    bool operator()(ListNode* a, ListNode* b) {
        return a->val > b->val; // min-heap by value
    }
};

ListNode* mergeKLists(vector<ListNode*>& lists) {
    priority_queue<ListNode*, vector<ListNode*>, Compare> minHeap;

    for (ListNode* node : lists) {
        if (node) minHeap.push(node);
    }

    ListNode dummy(0);
    ListNode* tail = &dummy;

    while (!minHeap.empty()) {
        ListNode* node = minHeap.top();
        minHeap.pop();

        tail->next = node;
        tail = node;

        if (node->next) minHeap.push(node->next);
    }
    return dummy.next;
}

int main() {
    ListNode* l1 = new ListNode(1);
    l1->next = new ListNode(4);
    l1->next->next = new ListNode(5);

    ListNode* l2 = new ListNode(1);
    l2->next = new ListNode(3);
    l2->next->next = new ListNode(4);

    ListNode* l3 = new ListNode(2);
    l3->next = new ListNode(6);

    vector<ListNode*> lists = {l1, l2, l3};
    ListNode* merged = mergeKLists(lists);

    for (ListNode* cur = merged; cur; cur = cur->next) {
        cout << cur->val << " ";
    }
    cout << endl;
    return 0;
}
\end{lstlisting}

\begin{complexitybox}
\textbf{Time Complexity.} The heap never holds more than $k$ nodes. Every
one of the $n$ total nodes (across all $k$ lists) is pushed once and
popped once, each costing $O(\log k)$, giving an overall time complexity
of
\[
  O(n \log k).
\]
\textbf{Space Complexity.} $O(k)$ for the heap (the output list reuses
the input nodes rather than allocating new ones).
\end{complexitybox}

\begin{takeawaybox}
``Merge $k$ sources, each individually sorted'' always reduces to ``track
one current candidate per source in a min-heap, pop the global minimum,
advance that source.'' This template applies identically whether the
sources are linked lists (this section), arrays, or streams, and is one of
the most directly reusable patterns in this chapter.
\end{takeawaybox}
